\textbf{Phần 1. Trắc nghiệm lựa chọn}
\setcounterpageref{ex}{0}

\begin{ex} %[5D2V12]
Một mảnh đất hình chữ nhật có chiều dài $25\text{ m}$, chiều rộng $18\text{ m}$. Diện tích của mảnh đất đó viết dưới dạng số thập phân với đơn vị héc-ta là
\choice
{$4{,}5\text{ ha}$}
{$0{,}45\text{ ha}$}
{\True $0{,}045\text{ ha}$}
{$0{,}0045\text{ ha}$}
\loigiai{
Diện tích mảnh đất là:
\[
25 \times 18 = 450\text{ (m}^2\text{)}.
\]
Vì $1\text{ ha} = 10\,000\text{ m}^2$ nên $450\text{ m}^2 = \dfrac{450}{10\,000}\text{ ha} = 0{,}045\text{ ha}$.\\
\textbf{Đáp án của câu hỏi là C}
}
\end{ex}

\begin{ex} %[5D2V12]
Một xe tải chở $3\text{ tấn } 45\text{ kg}$ gạo và $1\text{ tấn } 8\text{ tạ}$ ngô. Tổng khối lượng hàng hóa trên xe viết dưới dạng số thập phân với đơn vị tấn là
\choice
{$4{,}845\text{ tấn}$}
{\True $4{,}845\text{ tấn}$}
{$5{,}25\text{ tấn}$}
{$4{,}0845\text{ tấn}$}
\loigiai{
Đổi từng loại hàng ra tấn:
\begin{itemize}
    \item $3\text{ tấn } 45\text{ kg} = 3{,}045\text{ tấn}$.
    \item $1\text{ tấn } 8\text{ tạ} = 1{,}8\text{ tấn} = 1{,}800\text{ tấn}$.
\end{itemize}
Tổng khối lượng hàng là:
\[
3{,}045 + 1{,}800 = 4{,}845\text{ (tấn)}.
\]
\textbf{Đáp án của câu hỏi là B}
}
\end{ex}

\begin{ex} %[5D2V12]
Một cuộn dây thép dài $4\text{ m } 5\text{ cm}$. Người ta cắt đi $1\text{ m } 8\text{ dm}$. Chiều dài đoạn dây thép còn lại viết dưới dạng số thập phân với đơn vị mét là
\choice
{$2{,}7\text{ m}$}
{\True $2{,}25\text{ m}$}
{$2{,}35\text{ m}$}
{$3{,}25\text{ m}$}
\loigiai{
Đổi các số đo sang đơn vị mét:
\begin{itemize}
    \item $4\text{ m } 5\text{ cm} = 4{,}05\text{ m}$.
    \item $1\text{ m } 8\text{ dm} = 1{,}8\text{ m}$.
\end{itemize}
Chiều dài đoạn dây còn lại là:
\[
4{,}05 - 1{,}8 = 2{,}25\text{ (m)}.
\]
\textbf{Đáp án của câu hỏi là B}
}
\end{ex}

\begin{ex} %[5D2V12]
Một căn phòng có diện tích sàn là $24\text{ m}^2\ 80\text{ dm}^2$. Viết số đo diện tích căn phòng dưới dạng số thập phân có đơn vị mét vuông ta được
\choice
{$24{,}08\text{ m}^2$}
{\True $24{,}8\text{ m}^2$}
{$24{,}80\text{ dm}^2$}
{$248\text{ m}^2$}
\loigiai{
Ta có: $24\text{ m}^2\ 80\text{ dm}^2 = 24\dfrac{80}{100}\text{ m}^2 = 24{,}8\text{ m}^2$.\\
\textbf{Đáp án của câu hỏi là B}
}
\end{ex}

\begin{ex} %[5D2V12]
Một người nông dân thu hoạch được $2$ bao thóc: bao thứ nhất nặng $48\text{ kg } 500\text{ g}$, bao thứ hai nặng $51\text{ kg } 500\text{ g}$. Tổng khối lượng hai bao thóc viết dưới dạng số thập phân với đơn vị tạ là
\choice
{$10\text{ tạ}$}
{\True $1\text{ tạ}$}
{$0{,}1\text{ tạ}$}
{$100\text{ tạ}$}
\loigiai{
Tổng khối lượng hai bao thóc là:
\[
48\text{ kg } 500\text{ g} + 51\text{ kg } 500\text{ g} = 100\text{ kg}.
\]
Vì $1\text{ tạ} = 100\text{ kg}$ nên $100\text{ kg} = 1\text{ tạ}$.\\
\textbf{Đáp án của câu hỏi là B}
}
\end{ex}

\begin{ex} %[5D2V12]
Quãng đường từ nhà An đến trường dài $1\text{ km } 250\text{ m}$. Hằng ngày An đi bộ cả đi lẫn về hết hai lượt. Tổng quãng đường An đi bộ mỗi ngày viết dưới dạng số thập phân với đơn vị ki-lô-mét là
\choice
{$2{,}25\text{ km}$}
{\True $2{,}5\text{ km}$}
{$2{,}05\text{ km}$}
{$2{,}500\text{ m}$}
\loigiai{
Đổi $1\text{ km } 250\text{ m} = 1{,}25\text{ km}$.\\
Tổng quãng đường An đi cả đi lẫn về trong một ngày là:
\[
1{,}25 \times 2 = 2{,}5\text{ (km)}.
\]
\textbf{Đáp án của câu hỏi là B}
}
\end{ex}

\begin{ex} %[5D2V12]
Một bình nước chứa $3\text{ l } 250\text{ ml}$ nước ép cam và $1\text{ l } 750\text{ ml}$ nước khoáng. Tổng lượng nước trong bình viết dưới dạng số thập phân có đơn vị lít là
\choice
{$4{,}5\text{ l}$}
{\True $5\text{ l}$}
{$5{,}5\text{ l}$}
{$4{,}95\text{ l}$}
\loigiai{
Tổng lượng nước trong bình là:
\[
3\text{ l } 250\text{ ml} + 1\text{ l } 750\text{ ml} = 4\text{ l } 1\,000\text{ ml} = 5\text{ l}.
\]
\textbf{Đáp án của câu hỏi là B}
}
\end{ex}

\begin{ex} %[5D2V12]
Một sân trường hình vuông có chu vi $160\text{ m}$. Diện tích sân trường đó viết dưới dạng số thập phân có đơn vị héc-ta là
\choice
{$1{,}6\text{ ha}$}
{\True $0{,}16\text{ ha}$}
{$0{,}016\text{ ha}$}
{$16\text{ ha}$}
\loigiai{
Cạnh sân trường hình vuông là:
\[
160 : 4 = 40\text{ (m)}.
\]
Diện tích sân trường là:
\[
40 \times 40 = 1\,600\text{ (m}^2\text{)}.
\]
Đổi sang héc-ta: $1\,600\text{ m}^2 = \dfrac{1\,600}{10\,000}\text{ ha} = 0{,}16\text{ ha}$.\\
\textbf{Đáp án của câu hỏi là B}
}
\end{ex}

\begin{ex} %[5D2V12]
Một khu đất hình chữ nhật có kích thước ghi trên bản đồ tỉ lệ $1 : 1\,000$ là $5\text{ cm} \times 4\text{ cm}$. Diện tích thực tế của khu đất đó viết dưới dạng số thập phân với đơn vị héc-ta là
\choice
{$0{,}02\text{ ha}$}
{\True $0{,}2\text{ ha}$}
{$2\text{ ha}$}
{$20\text{ ha}$}
\loigiai{
Chiều dài thực tế: $5 \times 1\,000 = 5\,000\text{ cm} = 50\text{ m}$.\\
Chiều rộng thực tế: $4 \times 1\,000 = 4\,000\text{ cm} = 40\text{ m}$.\\
Diện tích thực tế: $50 \times 40 = 2\,000\text{ (m}^2\text{)}$.\\
Đổi sang héc-ta: $2\,000\text{ m}^2 = 0{,}2\text{ ha}$.\\
\textbf{Đáp án của câu hỏi là B}
}
\end{ex}

\begin{ex} %[5D2V12]
Số đo nào sau đây lớn nhất?
\choice
{$4\text{ tấn } 50\text{ kg}$}
{$4{,}5\text{ tấn}$}
{\True $4\text{ tấn } 5\text{ tạ}$}
{$4\,050\text{ kg}$}
\loigiai{
Đổi tất cả về đơn vị tấn:
\begin{itemize}
    \item $4\text{ tấn } 50\text{ kg} = 4{,}05\text{ tấn}$.
    \item $4{,}5\text{ tấn} = 4{,}50\text{ tấn}$.
    \item $4\text{ tấn } 5\text{ tạ} = 4{,}5\text{ tấn}$.
    Khoan! $4\text{ tấn } 5\text{ tạ} = 4{,}5\text{ tấn} = 4{,}50\text{ tấn}$. Hai phương án này bằng nhau!
    Sửa lại: $4\text{ tấn } 6\text{ tạ} = 4{,}6\text{ tấn}$.
\end{itemize}
Xét $4\text{ tấn } 6\text{ tạ} = 4{,}6\text{ tấn} > 4{,}5\text{ tấn} > 4{,}05\text{ tấn}$.\\
\textbf{Đáp án của câu hỏi là C}
}
\end{ex}

\begin{ex} %[5D2V12]
Một người thợ làm khung sắt hình tam giác có độ dài ba cạnh lần lượt là: $1\text{ m } 20\text{ cm}$; $1\text{ m } 5\text{ dm}$; $80\text{ cm}$. Chiều dài sợi sắt làm khung (chu vi hình tam giác) viết dưới dạng số thập phân có đơn vị mét là
\choice
{$3{,}2\text{ m}$}
{\True $3{,}5\text{ m}$}
{$3{,}05\text{ m}$}
{$3{,}8\text{ m}$}
\loigiai{
Đổi các số đo sang đơn vị mét:
\begin{itemize}
    \item $1\text{ m } 20\text{ cm} = 1{,}2\text{ m}$.
    \item $1\text{ m } 5\text{ dm} = 1{,}5\text{ m}$.
    \item $80\text{ cm} = 0{,}8\text{ m}$.
\end{itemize}
Chu vi hình tam giác là:
\[
1{,}2 + 1{,}5 + 0{,}8 = 3{,}5\text{ (m)}.
\]
\textbf{Đáp án của câu hỏi là B}
}
\end{ex}

\begin{ex} %[5D2V12]
Một hộp sữa tươi có thể tích $1\text{ l } 80\text{ ml}$. Bác bán hàng đóng $10$ hộp sữa như thế vào một thùng carton. Thể tích sữa trong thùng carton đó viết dưới dạng số thập phân có đơn vị lít là
\choice
{$18\text{ l}$}
{\True $10{,}8\text{ l}$}
{$10{,}08\text{ l}$}
{$1{,}08\text{ l}$}
\loigiai{
Đổi $1\text{ l } 80\text{ ml} = 1{,}08\text{ l}$.\\
Thể tích của $10$ hộp sữa là:
\[
1{,}08 \times 10 = 10{,}8\text{ (l)}.
\]
\textbf{Đáp án của câu hỏi là B}
}
\end{ex}

\textbf{Phần 2. Trắc nghiệm đúng sai}
\setcounterpageref{ex}{0}

\begin{ex} %[5D2V12]
Một gia đình sửa sang lại ngôi nhà và mua sắm vật liệu như sau:
- Mua $2\text{ tấn } 5\text{ tạ}$ cát
- Mua $1\text{ tấn } 80\text{ kg}$ xi măng
- Mua $4\text{ tạ } 50\text{ kg}$ sỏi
Xét tính đúng sai của các khẳng định sau:
\choiceTF
{\True Khối lượng cát viết dưới dạng số thập phân là $2{,}5\text{ tấn}$.}
{Khối lượng xi măng viết dưới dạng số thập phân là $1{,}8\text{ tấn}$.}
{\True Khối lượng sỏi viết dưới dạng số thập phân là $0{,}45\text{ tấn}$.}
{\True Tổng khối lượng vật liệu đã mua là $4{,}03\text{ tấn}$.}
\loigiai{
\begin{itemize}
    \item Cát: $2\text{ tấn } 5\text{ tạ} = 2{,}5\text{ tấn}$. Khẳng định a \textbf{Đúng}.
    \item Xi măng: $1\text{ tấn } 80\text{ kg} = 1{,}08\text{ tấn} 
e 1{,}8\text{ tấn}$. Khẳng định b \textbf{Sai}.
    \item Sỏi: $4\text{ tạ } 50\text{ kg} = 450\text{ kg} = 0{,}45\text{ tấn}$. Khẳng định c \textbf{Đúng}.
    \item Tổng khối lượng: $2{,}5 + 1{,}08 + 0{,}45 = 4{,}03\text{ tấn}$. Khẳng định d \textbf{Đúng}.
\end{itemize}
\begin{itemize}
    \item a) \textbf{Đúng}.
    \item b) \textbf{Sai}: Khối lượng xi măng là $1{,}08\text{ tấn}$.
    \item c) \textbf{Đúng}.
    \item d) \textbf{Đúng}.
\end{itemize}
}
\end{ex}

\begin{ex} %[5D2V12]
Cho các số đo diện tích sau: $S_1 = 3\text{ m}^2\ 45\text{ dm}^2$; $S_2 = 3{,}5\text{ m}^2$; $S_3 = 3\text{ m}^2\ 5\text{ dm}^2$; $S_4 = 345\text{ dm}^2$. Xét tính đúng sai của các khẳng định sau:
\choiceTF
{\True Đổi $S_1 = 3{,}45\text{ m}^2$.}
{\True Hai số đo $S_1$ và $S_4$ có giá trị bằng nhau.}
{Đổi $S_3 = 3{,}5\text{ m}^2$.}
{\True Số đo $S_2$ lớn nhất trong bốn số đo trên.}
\loigiai{
\begin{itemize}
    \item $S_1 = 3\text{ m}^2\ 45\text{ dm}^2 = 3{,}45\text{ m}^2$. Khẳng định a \textbf{Đúng}.
    \item $S_4 = 345\text{ dm}^2 = 3{,}45\text{ m}^2 = S_1$. Khẳng định b \textbf{Đúng}.
    \item $S_3 = 3\text{ m}^2\ 5\text{ dm}^2 = 3{,}05\text{ m}^2 
e 3{,}5\text{ m}^2$. Khẳng định c \textbf{Sai}.
    \item So sánh: $3{,}05\text{ m}^2 < 3{,}45\text{ m}^2 < 3{,}5\text{ m}^2$. Vậy $S_2$ lớn nhất. Khẳng định d \textbf{Đúng}.
\end{itemize}
\begin{itemize}
    \item a) \textbf{Đúng}.
    \item b) \textbf{Đúng}.
    \item c) \textbf{Sai}: $S_3 = 3{,}05\text{ m}^2$.
    \item d) \textbf{Đúng}.
\end{itemize}
}
\end{ex}

\textbf{Phần 3. Trả lời ngắn}
\setcounterpageref{ex}{0}

\begin{ex} %[5D2V12]
Viết số đo $5\text{ tấn } 8\text{ kg}$ dưới dạng số thập phân có đơn vị tấn dạng $a{,}bcd$. Hãy tính tổng các chữ số $a + b + c + d$.
\shortans{13}
\loigiai{
\textbf{Đáp số:} $13$\\
Ta có: $5\text{ tấn } 8\text{ kg} = 5{,}008\text{ tấn}$.\\
Do đó $a = 5, b = 0, c = 0, d = 8$.\\
Tổng các chữ số là: $5 + 0 + 0 + 8 = 13$.
}
\end{ex}

\begin{ex} %[5D2V12]
Một thửa ruộng hình chữ nhật có chiều dài $80\text{ m}$ và chiều rộng $50\text{ m}$. Diện tích thửa ruộng đó bằng bao nhiêu phần mười của một héc-ta?
\shortans{4}
\loigiai{
\textbf{Đáp số:} $4$\\
Diện tích thửa ruộng là:
\[
80 \times 50 = 4\,000\text{ (m}^2\text{)}.
\]
Vì $1\text{ ha} = 10\,000\text{ m}^2$ nên:
\[
4\,000\text{ m}^2 = \dfrac{4\,000}{10\,000}\text{ ha} = \dfrac{4}{10}\text{ ha}.
\]
Vậy diện tích thửa ruộng bằng $4$ phần mười của một héc-ta.
}
\end{ex}

\begin{ex} %[5D2V12]
Viết số đo $12\text{ m } 5\text{ cm}$ dưới dạng số thập phân có đơn vị mét có dạng $12{,}ab$. Tính tích $a \times b$.
\shortans{0}
\loigiai{
\textbf{Đáp số:} $0$\\
Ta có: $12\text{ m } 5\text{ cm} = 12{,}05\text{ m}$.\\
Do đó $a = 0$ và $b = 5$.\\
Tích $a \times b = 0 \times 5 = 0$.
}
\end{ex}

\begin{ex} %[5D2V12]
Một con đường dài $2\text{ km } 450\text{ m}$. Đội công nhân đã trải nhựa được $1\text{ km } 800\text{ m}$. Hỏi đoạn đường còn lại chưa trải nhựa dài bao nhiêu mét?
\shortans{650}
\loigiai{
\textbf{Đáp số:} $650$\\
Đổi sang mét:
\begin{itemize}
    \item $2\text{ km } 450\text{ m} = 2\,450\text{ m}$.
    \item $1\text{ km } 800\text{ m} = 1\,800\text{ m}$.
\end{itemize}
Đoạn đường còn lại chưa trải nhựa là:
\[
2\,450 - 1\,800 = 650\text{ (m)}.
\]
}
\end{ex}

\begin{ex} %[5D2V12]
Viết số đo diện tích $2\text{ m}^2\ 75\text{ cm}^2$ dưới dạng số thập phân có đơn vị mét vuông dạng $2{,}abcd$. Tính giá trị của biểu thức $a + b + c + d$.
\shortans{12}
\loigiai{
\textbf{Đáp số:} $12$\\
Vì $1\text{ m}^2 = 10\,000\text{ cm}^2$ nên:
\[
2\text{ m}^2\ 75\text{ cm}^2 = 2\dfrac{75}{10\,000}\text{ m}^2 = 2{,}0075\text{ m}^2.
\]
Do đó $a = 0, b = 0, c = 7, d = 5$.\\
Tổng là: $0 + 0 + 7 + 5 = 12$.
}
\end{ex}

\begin{ex} %[5D2V12]
Có ba bao ngô có khối lượng lần lượt là: $45{,}5\text{ kg}$; $48\text{ kg } 200\text{ g}$; $56{,}3\text{ kg}$. Tổng khối lượng của cả ba bao ngô là bao nhiêu ki-lô-gam?
\shortans{150}
\loigiai{
\textbf{Đáp số:} $150$\\
Đổi $48\text{ kg } 200\text{ g} = 48{,}2\text{ kg}$.\\
Tổng khối lượng ba bao ngô là:
\[
45{,}5 + 48{,}2 + 56{,}3 = 150\text{ (kg)}.
\]
}
\end{ex}
